%% abstract.tex — P7 Quantum Molecular Encoding (REFOCUSED)
%% Status: STUB — pending canonical Phase D results
%% Last updated: 06 October 2026 (refocus checkpoint)
%% All \PENDING{} markers must be filled from DAR before submission.

African natural products (ANPs) present a persistent representation challenge
for computational antimalarial screening: ECFP4 fingerprints collapse
stereoisomers into identical bit vectors, and graph neural network message-passing
cannot distinguish cis/trans double bonds or polycyclic ring geometries across
novel scaffolds — limitations documented across Projects 1, 3, and 5 of this
series. Quantum molecular structure encoding (QMSE), introduced by Boy et al.
(\textit{Mach.\ Learn.: Sci.\ Technol.}\ \textbf{6}, 045076, 2025), encodes
bond orders, Z/E stereochemistry, and R/S chirality directly as one- and
two-qubit rotation angles in parameterised quantum circuits — bypassing
fingerprint hashing entirely. We investigate whether a hybrid quantum-classical
model coupling QMSE circuits (PennyLane \texttt{default.qubit}) to a TensorFlow
dense head, trained end-to-end via automatic differentiation through both
components, can overcome these structural encoding limitations for antimalarial
activity prediction on African natural product-enriched datasets.

We encode \PENDING{n\_valid} molecules from the \num{19849}-compound P3 benchmark
using \PENDING{BondOrderMatrix / CoulombMatrix} and evaluate the hybrid QNN under
stratified 5-fold cross-validation and 100\% novel-scaffold holdout, comparing
against ECFP4-RBF SVM (\text{AUC} = 0.9475, P3 canonical), the P3 quantum-inspired
kernel (\text{AUC} = 0.8385), and the prior ECFP4-VQC exploratory baseline
(\text{AUC} = 0.8474, P7 Phase 4 archived). Hyperparameter optimisation via
Optuna TPE (\PENDING{n\_trials} trials, 3-fold inner CV) identifies the optimal
configuration across encoding type, circuit depth, entangling gate, and classical
head architecture. The central finding is \PENDING{Scenario~A/B/C}: \PENDING{one-line
result statement and canonical AUC}. Gradient diagnostics confirm \PENDING{barren
plateau status}. \PENDING{Closing sentence: implication for quantum advantage,
limitations, and next step toward hardware validation.}
